% Based on Li2_性质与特殊值汇总.pdf (2026-10-21); corrected and expanded proofs.
% Same mathematical content and equation numbers as the companion Markdown.
% Build from repository root, run twice:
% xelatex -no-shell-escape -halt-on-error -interaction=nonstopmode -output-directory=outputs/dilogarithm print/dilogarithm-properties-special-values.tex
% Approved layout: A4, 23 mm margins, 11 pt body, 1.12 baseline stretch,
% 3 pt paragraph skip; Fandol Song and original Computer Modern mathematics.
\documentclass[11pt,a4paper]{article}
\usepackage[margin=23mm,headheight=14pt,headsep=8.5mm,footskip=10mm]{geometry}
\usepackage{amsmath,amssymb,amsthm}
\usepackage[no-math]{fontspec}
\usepackage{xeCJK}
\setmainfont{Latin Modern Roman}
\setCJKmainfont{FandolSong-Regular.otf}[Script=CustomDefault,Path=assets/fonts/fandol/,BoldFont=FandolSong-Bold.otf,ItalicFont=FandolSong-Regular.otf,BoldItalicFont=FandolSong-Bold.otf]
\setCJKsansfont{FandolSong-Regular.otf}[Script=CustomDefault,Path=assets/fonts/fandol/,BoldFont=FandolSong-Bold.otf]
\xeCJKsetup{PunctStyle=quanjiao,CheckSingle=true}
\usepackage{setspace}
\setstretch{1.12}
\setlength{\parindent}{2em}
\setlength{\parskip}{3pt}
\usepackage{titlesec}
\titleformat{\section}{\centering\normalfont\Large\bfseries}{\thesection.}{0.6em}{}
\titleformat{\subsection}{\normalfont\large\bfseries}{\thesubsection}{0.6em}{}
\titlespacing*{\section}{0pt}{1.8\baselineskip}{1.0\baselineskip}
\titlespacing*{\subsection}{0pt}{1.1\baselineskip}{0.5\baselineskip}
\usepackage{enumitem}
\setlist[enumerate]{leftmargin=2em,labelsep=0.5em,itemsep=3pt,parsep=0pt,topsep=6pt}
\providecommand{\tightlist}{\setlength{\itemsep}{3pt}\setlength{\parskip}{0pt}}
\usepackage{needspace}
\newenvironment{statementblock}{\par\Needspace{9\baselineskip}\noindent}{\par\nopagebreak[4]}
\newenvironment{proofblock}{\par\Needspace{4\baselineskip}\noindent}{\par}
\usepackage{fancyhdr}
\pagestyle{fancy}
\fancyhf{}
\fancyhead[L]{\small 数学手记}
\fancyhead[R]{\small 二对数函数}
\fancyfoot[C]{\thepage}
\renewcommand{\headrulewidth}{0.4pt}
\usepackage[hidelinks,unicode,bookmarksnumbered=true,linktoc=all]{hyperref}
\usepackage{bookmark}
\bookmarksetup{depth=2}
\hypersetup{pdftitle={二对数函数的性质与特殊值计算},pdfauthor={Vincent Arqain},pdfsubject={原三页讲义的校订版：补全实轴延拓、证明与误差界}}
\renewcommand{\contentsname}{目录}
\setcounter{tocdepth}{1}
\setcounter{secnumdepth}{2}
\setcounter{secnumdepth}{2}
\widowpenalty=10000
\clubpenalty=10000
\displaywidowpenalty=10000
\predisplaypenalty=10000
\interdisplaylinepenalty=10000
\raggedbottom
\begin{document}
\begin{center}
{\LARGE\bfseries 二对数函数的性质与特殊值计算\par}
\vspace{8pt}
{\normalsize 定义、函数方程与典型应用\par}
\end{center}
\vspace{10pt}
本文由《Li2\_性质与特殊值汇总》校订整理，保留原有特殊值计算与方法总结，补全实轴延拓、函数方程的证明和渐近余项说明。所有变量均为实数，所有对数均为自然对数。本文讨论的函数定义在 \((-\infty,1]\) 上，不涉及 \(x>1\) 或复平面上的分支选择。

与本文配套的 \href{https://vincent-arqain.github.io/math-notes/notes/890e1516e64e21a2fb2248b4/}{Euler 型级数讲义} 系统讨论换序、调和数生成函数与多重 Zeta 值；本文把其中最常用的二对数函数单独展开。本版与配套网页使用相同的公式编号。

\hypertarget{ux5b9aux4e49ux4e0eux57faux672cux6027ux8d28}{%
\section{定义与基本性质}\label{ux5b9aux4e49ux4e0eux57faux672cux6027ux8d28}}

\hypertarget{ux5e42ux7ea7ux6570ux7684ux6536ux655bux8303ux56f4}{%
\subsection{幂级数的收敛范围}\label{ux5e42ux7ea7ux6570ux7684ux6536ux655bux8303ux56f4}}

对 \(|x|\leqslant1\)，定义

\[
\operatorname{Li}_2(x)=\sum_{n=1}^{\infty}\frac{x^n}{n^2}.
\]

因为 \(|x^n/n^2|\leqslant1/n^2\)，该级数在 \([-1,1]\) 上一致绝对收敛，故其和函数连续。若 \(|x|>1\)，通项不趋于零，级数发散；特别是不能用这个级数直接定义 \(x<-1\) 时的值。

在 \(|x|<1\) 内，对 \(-\log(1-x)=\sum_{n\geqslant1}x^n/n\) 除以 \(x\)，再在从 \(0\) 到 \(x\) 的紧区间上逐项积分，得到积分表示。这里的一致收敛只用在收敛区间内部。

\hypertarget{ux79efux5206ux5ef6ux62d3ux4e0eux7aefux70b9}{%
\subsection{积分延拓与端点}\label{ux79efux5206ux5ef6ux62d3ux4e0eux7aefux70b9}}

对任意 \(x<1\)，用下式将函数延拓到整个负实轴：

\[
\operatorname{Li}_2(x)=-\int_0^x\frac{\log(1-t)}t\,\mathrm dt.
\]

积分核 \(-\log(1-t)/t\) 在 \(t=0\) 处的连续延拓值为 \(1\)，所以跨过零点没有奇性；当 \(x<0\) 时，积分按通常的有向积分理解。该定义与幂级数在 \((-1,1)\) 内一致，在 \(x=-1\) 处由连续性一致。对于 \(x=-a<0\)，也可写成

\[
\operatorname{Li}_2(-a)=-\int_0^a\frac{\log(1+u)}u\,\mathrm du,
\qquad a>0.
\]

当 \(x\uparrow1\) 时，积分中的对数奇性可积：在 \([1/2,1)\) 上，\(1/t\leqslant2\)，而 \(-\log(1-t)\) 可积。因此积分定义连续延伸到 \(x=1\)，并与一致收敛的级数定义相同。于是

\[
\operatorname{Li}_2(0)=0,\qquad
\operatorname{Li}_2(1)=\zeta(2)=\frac{\pi^2}{6},\qquad
\operatorname{Li}_2(-1)=-\frac{\pi^2}{12}.
\]

最后一个等式由绝对收敛级数的奇偶分组得到：\(\sum(-1)^n/n^2=2\sum(2n)^{-2}-\zeta(2)=-\zeta(2)/2\)。所用的 Basel 求和值 \(\zeta(2)=\pi^2/6\)，可见 \href{https://vincent-arqain.github.io/math-notes/notes/890e1516e64e21a2fb2248b4/\#计算二次特殊值}{Euler 讲义的附录证明}。

\hypertarget{ux5bfcux6570ux5355ux8c03ux6027ux4e0eux51f8ux6027}{%
\subsection{导数、单调性与凸性}\label{ux5bfcux6570ux5355ux8c03ux6027ux4e0eux51f8ux6027}}

由微积分基本定理，对 \(x<1\)、\(x\ne0\)，

\[
\operatorname{Li}_2'(x)=-\frac{\log(1-x)}x,
\qquad
\operatorname{Li}_2''(x)=\frac{x/(1-x)+\log(1-x)}{x^2}.
\]

在零点由幂级数得到 \(\operatorname{Li}_2'(0)=1\)、\(\operatorname{Li}_2''(0)=1/2\)。更统一的表达式是

\[
\operatorname{Li}_2'(x)=\int_0^1\frac{\mathrm du}{1-xu}>0,
\qquad
\operatorname{Li}_2''(x)=\int_0^1\frac{u\,\mathrm du}{(1-xu)^2}>0,
\qquad x<1.
\]

第一个积分可直接计算；第二个由对第一个积分求导得到。在任一固定的 \(x<1\) 的小闭邻域内，分母一致远离零，因而该求导合法。故函数在 \((-\infty,1)\) 上严格递增、严格凸。导数公式不包含端点 \(x=1\)；事实上 \(\operatorname{Li}_2'(x)\to+\infty\)，差商的平均值表达式也说明 \(x=1\) 处没有有限左导数。

\hypertarget{ux56dbux4e2aux57faux672cux51fdux6570ux65b9ux7a0b}{%
\section{四个基本函数方程}\label{ux56dbux4e2aux57faux672cux51fdux6570ux65b9ux7a0b}}

\hypertarget{euler-ux53cdux5c04ux516cux5f0f}{%
\subsection{Euler 反射公式}\label{euler-ux53cdux5c04ux516cux5f0f}}

对 \(0<x<1\)，有

\begin{equation}
\operatorname{Li}_2(x)+\operatorname{Li}_2(1-x)
=\frac{\pi^2}{6}-\log x\log(1-x).
\tag{1}
\end{equation}

证明：令 \(F(x)=\operatorname{Li}_2(x)+\operatorname{Li}_2(1-x)+\log x\log(1-x)\)，则

\[
F'(x)=-\frac{\log(1-x)}x+\frac{\log x}{1-x}
+\frac{\log(1-x)}x-\frac{\log x}{1-x}=0.
\]

由端点连续性及 \(\log x\log(1-x)=O(x|\log x|)\to0\)，令 \(x\downarrow0\) 即知常数为 \(\pi^2/6\)。式 (1) 在端点只按极限解释，不能把 \(\log0\) 当作实数代入。

\hypertarget{landen-ux53d8ux6362ux516cux5f0f}{%
\subsection{Landen 变换公式}\label{landen-ux53d8ux6362ux516cux5f0f}}

对所有 \(x<1\)，有

\begin{equation}
\operatorname{Li}_2(x)+\operatorname{Li}_2\!\left(\frac{x}{x-1}\right)
=-\frac12\log^2(1-x).
\tag{2}
\end{equation}

证明：记 \(y=x/(x-1)\)。因为 \(1-y=1/(1-x)>0\)，所以 \(y<1\)，两个函数值均在积分定义域内。令左端加上 \(\tfrac12\log^2(1-x)\) 为 \(G(x)\)。对 \(x\ne0\)，链式法则给出

\[
G'(x)=-\frac{\log(1-x)}x
+\frac{\log(1-x)}{x(1-x)}
-\frac{\log(1-x)}{1-x}=0.
\]

因此 \(G\) 在 \((-\infty,0)\) 和 \((0,1)\) 上分别为常数。它在零点连续，且 \(G(0)=0\)，故两个常数均为零。式 (2) 包含 \(x=0\)。

\hypertarget{ux500dux89d2ux516cux5f0f}{%
\subsection{倍角公式}\label{ux500dux89d2ux516cux5f0f}}

对 \(|x|\leqslant1\)，有

\begin{equation}
\operatorname{Li}_2(x)+\operatorname{Li}_2(-x)
=\frac12\operatorname{Li}_2(x^2).
\tag{3}
\end{equation}

证明：绝对收敛允许逐项相加并取偶数项，故

\[
\sum_{n=1}^{\infty}\frac{x^n+(-x)^n}{n^2}
=2\sum_{k=1}^{\infty}\frac{x^{2k}}{(2k)^2}
=\frac12\operatorname{Li}_2(x^2).
\]

这个证明包含 \(x=\pm1\)。本文不把式 (3) 延伸到 \(|x|>1\)，因为那时 \(x^2>1\) 已超出本文的实值定义域。

\hypertarget{ux8d1fux5b9eux8f74ux53cdux6f14ux516cux5f0f}{%
\subsection{负实轴反演公式}\label{ux8d1fux5b9eux8f74ux53cdux6f14ux516cux5f0f}}

对 \(a>0\)，有

\begin{equation}
\operatorname{Li}_2(-a)+\operatorname{Li}_2(-1/a)
=-\frac{\pi^2}{6}-\frac12\log^2a.
\tag{4}
\end{equation}

证明：令左端加上 \(\tfrac12\log^2a\) 为 \(J(a)\)，则

\[
J'(a)=\frac{-\log(1+a)+\log(1+1/a)+\log a}{a}=0,
\]

其中使用 \(1+1/a=(1+a)/a\) 和实对数的乘法公式。取 \(a=1\)，得到 \(J(1)=2\operatorname{Li}_2(-1)=-\pi^2/6\)。这里两个自变量始终为负数，任意大小的 \(a>0\) 均可使用积分延拓。

\hypertarget{ux7279ux6b8aux503cux53caux5176ux8ba1ux7b97ux65b9ux6cd5}{%
\section{特殊值及其计算方法}\label{ux7279ux6b8aux503cux53caux5176ux8ba1ux7b97ux65b9ux6cd5}}

\hypertarget{ux534aux70b9ux7684ux7279ux6b8aux503c}{%
\subsection{半点的特殊值}\label{ux534aux70b9ux7684ux7279ux6b8aux503c}}

在式 (1) 中取 \(x=1/2\)，即得

\begin{equation}
\operatorname{Li}_2\!\left(\frac12\right)
=\frac{\pi^2}{12}-\frac12\log^22.
\tag{5}
\end{equation}

这一步的关键是反射后自变量不变，从而原来的两个函数值合并为一个未知量。

\hypertarget{ux9ec4ux91d1ux5206ux5272ux6bd4ux5bf9ux5e94ux7684ux7279ux6b8aux503c}{%
\subsection{黄金分割比对应的特殊值}\label{ux9ec4ux91d1ux5206ux5272ux6bd4ux5bf9ux5e94ux7684ux7279ux6b8aux503c}}

令 \(\varphi=(1+\sqrt5)/2\)、\(t=\varphi^{-1}\)、\(L=\log\varphi\)。有 \(1-t=t^2\)、\(\log t=-L\)。记

\[
A=\operatorname{Li}_2(t),\quad B=\operatorname{Li}_2(t^2),\quad
C=\operatorname{Li}_2(-t),\quad D=\operatorname{Li}_2(-\varphi).
\]

\Needspace{9\baselineskip}
由于 \(0<t<1\)，式 (1)、(3)、(2)、(4) 分别给出

\[
\begin{aligned}
A+B&=\frac{\pi^2}{6}-2L^2,\\
A+C&=\frac12B,\\
A+D&=-2L^2,\\
C+D&=-\frac{\pi^2}{6}-\frac12L^2.
\end{aligned}
\]

第三式使用 \(t/(t-1)=-1/t=-\varphi\)；第四式取 \(a=\varphi\)。消去 \(C,D\) 得 \(4A-B=\pi^2/3-3L^2\)。与第一式联立，得到

\[
\begin{aligned}
\operatorname{Li}_2\!\left(\frac{\sqrt5-1}{2}\right)
&=\frac{\pi^2}{10}-\log^2\varphi,\\
\operatorname{Li}_2\!\left(\frac{3-\sqrt5}{2}\right)
&=\frac{\pi^2}{15}-\log^2\varphi.
\end{aligned}
\]

再由 \(C=B/2-A\)、\(D=-2L^2-A\)，得

\[
\operatorname{Li}_2(-\varphi^{-1})=-\frac{\pi^2}{15}+\frac12\log^2\varphi,
\qquad
\operatorname{Li}_2(-\varphi)=-\frac{\pi^2}{10}-\log^2\varphi.
\]

这里并没有把 \(-\varphi<-1\) 代入发散的幂级数；该值来自积分延拓和已经证明的函数方程。

\hypertarget{ux6709ux7406ux81eaux53d8ux91cfux4e4bux95f4ux7684ux5173ux7cfb}{%
\subsection{有理自变量之间的关系}\label{ux6709ux7406ux81eaux53d8ux91cfux4e4bux95f4ux7684ux5173ux7cfb}}

对一般有理自变量，本讲义不保证其值能单独化成 \(\pi^2\) 与对数的简单组合。但可以证明多个值之间的恒等式，例如

\begin{equation}
\operatorname{Li}_2\!\left(\frac14\right)
+2\operatorname{Li}_2\!\left(\frac13\right)
=\frac{\pi^2}{6}-\log^22-\log^2\!\frac32.
\tag{6}
\end{equation}

证明：式 (3) 取 \(x=1/2\)，式 (2) 取 \(x=-1/2\)，分别得到

\[
\begin{aligned}
\operatorname{Li}_2(1/2)+\operatorname{Li}_2(-1/2)
&=\frac12\operatorname{Li}_2(1/4),\\
\operatorname{Li}_2(-1/2)+\operatorname{Li}_2(1/3)
&=-\frac12\log^2\!\frac32.
\end{aligned}
\]

消去 \(\operatorname{Li}_2(-1/2)\) 并代入式 (5)，即得式 (6)。同时保留有用的关系

\[
\operatorname{Li}_2(-1/2)=-\operatorname{Li}_2(1/3)-\frac12\log^2\!\frac32.
\]

这条具体恒等式只需倍角与 Landen 公式，无需引入 Abel 五项关系。

\hypertarget{euler-ux578bux7ea7ux6570ux4e0eux5bf9ux6570ux79efux5206}{%
\section{Euler 型级数与对数积分}\label{euler-ux578bux7ea7ux6570ux4e0eux5bf9ux6570ux79efux5206}}

\hypertarget{ux8c03ux548cux6570ux751fux6210ux51fdux6570ux4e0eux6c42ux548c}{%
\subsection{调和数生成函数与求和}\label{ux8c03ux548cux6570ux751fux6210ux51fdux6570ux4e0eux6c42ux548c}}

记 \(H_n=\sum_{k=1}^n1/k\)。对 \(|x|<1\)，双重级数绝对收敛，因为令 \(r=|x|\) 有

\[
\sum_{n=1}^{\infty}\sum_{k=1}^n\frac{r^n}{k}
=\frac1{1-r}\sum_{k=1}^{\infty}\frac{r^k}{k}<\infty.
\]

右端可先由非负项换序得到，再据其有限性判断绝对收敛。因此对一般实数 \(|x|<1\)，换序合法，且

\[
\sum_{n=1}^{\infty}H_nx^n=-\frac{\log(1-x)}{1-x}.
\]

在收敛区间内部，除以 \(x\) 后的级数在零点也连续，并可逐项积分。利用 \(1/[x(1-x)]=1/x+1/(1-x)\)，得到

\begin{equation}
\sum_{n=1}^{\infty}\frac{H_n}{n}x^n
=\operatorname{Li}_2(x)+\frac12\log^2(1-x),
\qquad |x|<1.
\tag{7}
\end{equation}

在式 (7) 中取 \(x=1/2\)，再用式 (5)，两个对数平方项抵消，故

\[
\sum_{n=1}^{\infty}\frac{H_n}{n\,2^n}=\frac{\pi^2}{12}.
\]

注意分母是 \(n\,2^n\)，并非 \(n^2\)。若在式 (7) 两边再除以 \(x\)，并在 \((0,1)\) 上积分，非负项逐项积分给出

\[
\sum_{n=1}^{\infty}\frac{H_n}{n^2}
=\zeta(3)+\frac12\int_0^1\frac{\log^2(1-x)}x\,\mathrm dx
=2\zeta(3).
\]

第一项由 \(\operatorname{Li}_2(x)/x\) 的非负幂级数积分得到。对第二项作 \(u=1-x\)，再以非负几何级数展开 \(1/(1-u)\)，使用

\[
\int_0^1u^j\log^2u\,\mathrm du=\frac{2}{(j+1)^3}
\]

即可得到积分等于 \(2\zeta(3)\)；此矩积分由 \(u=e^{-v}\) 和两次分部积分得出。上述步骤同时证明积分有限。更系统的推导与例题见 \href{https://vincent-arqain.github.io/math-notes/notes/890e1516e64e21a2fb2248b4/\#生成函数与除以指标}{Euler 讲义的生成函数章节} 及 \href{https://vincent-arqain.github.io/math-notes/notes/890e1516e64e21a2fb2248b4/\#例-4严格计算三次权的基本-euler-和}{三次权基本 Euler 和}。式 (7) 的端点不能不经收敛性分析就直接代入。

\Needspace{11\baselineskip}
\hypertarget{ux4e24ux4e2aux5178ux578bux5bf9ux6570ux79efux5206}{%
\subsection{两个典型对数积分}\label{ux4e24ux4e2aux5178ux578bux5bf9ux6570ux79efux5206}}

解：由积分定义和已求出的特殊值，

\[
\begin{aligned}
\int_0^1\frac{\log(1+x)}x\,\mathrm dx
&=-\operatorname{Li}_2(-1)=\frac{\pi^2}{12},\\
\int_0^{1/2}\frac{\log(1-x)}x\,\mathrm dx
&=-\operatorname{Li}_2(1/2)=-\frac{\pi^2}{12}+\frac12\log^22.
\end{aligned}
\]

零点处两个被积函数的极限分别为 \(1\) 和 \(-1\)，故不存在不可积奇性。遇到``对数除以自变量''的积分，先检查是否可以直接识别为二对数函数的差值。

\hypertarget{ux6570ux503cux8ba1ux7b97ux4e0eux6e10ux8fd1ux5c55ux5f00}{%
\section{数值计算与渐近展开}\label{ux6570ux503cux8ba1ux7b97ux4e0eux6e10ux8fd1ux5c55ux5f00}}

\hypertarget{ux622aux65adux8befux5deeux4e0eux81eaux53d8ux91cfux7ea6ux5316}{%
\subsection{截断误差与自变量约化}\label{ux622aux65adux8befux5deeux4e0eux81eaux53d8ux91cfux7ea6ux5316}}

对 \(0\leqslant r<1\) 及整数 \(N\geqslant0\)，有

\begin{equation}
0\leqslant\operatorname{Li}_2(r)-\sum_{n=1}^N\frac{r^n}{n^2}
\leqslant\frac{r^{N+1}}{(N+1)^2(1-r)}.
\tag{8}
\end{equation}

证明：余项的每个分母均不小于 \((N+1)^2\)，再求几何尾和即可。对 \(|x|<1\)，绝对误差也被式 (8) 中取 \(r=|x|\) 的右端控制；端点 \(r=1\) 不使用这个界。

实际计算可按以下步骤将级数参数约化到 \([0,1/2]\)。

\begin{enumerate}
\def\labelenumi{\arabic{enumi}.}
\tightlist
\item
  \(0\leqslant x\leqslant1/2\)：直接用级数和式 (8)。
\item
  \(1/2<x<1\)：用式 (1) 化为 \(1-x\)。
\item
  \(-1\leqslant x<0\)：用式 (2) 化为 \(x/(x-1)\in(0,1/2]\)。
\item
  \(x<-1\)：令 \(a=-x>1\)，先用式 (4) 化为 \(-1/a\in(-1,0)\)，再用上一步。
\item
  \(x=1\)：直接使用 \(\pi^2/6\)。
\end{enumerate}

这个余项界控制的是截断误差；用浮点数计算对数、平方和相减时，还应保留足够精度，不能将截断界误当作全部舍入误差的界。

\hypertarget{ux96f6ux70b9ux9644ux8fd1}{%
\subsection{零点附近}\label{ux96f6ux70b9ux9644ux8fd1}}

原级数直接给出

\[
\operatorname{Li}_2(x)=x+\frac{x^2}{4}+\frac{x^3}{9}+O(x^4),
\qquad x\to0.
\]

由式 (8) 的绝对值版本，对 \(|x|\leqslant1/2\)，前三项之后的余项绝对值不超过 \(|x|^4/[16(1-|x|)]\leqslant|x|^4/8\)，故这里的大 \(O\) 是双侧且一致的。

\Needspace{10\baselineskip}
\hypertarget{ux6b63ux7aefux70b9ux9644ux8fd1}{%
\subsection{正端点附近}\label{ux6b63ux7aefux70b9ux9644ux8fd1}}

对 \(\varepsilon\downarrow0\)，式 (1) 给出

\[
\operatorname{Li}_2(1-\varepsilon)
=\frac{\pi^2}{6}+\varepsilon\log\varepsilon-\varepsilon
+O\!\left(\varepsilon^2|\log\varepsilon|\right).
\]

证明：在 \(0<\varepsilon\leqslant1/2\) 上，

\[
|\operatorname{Li}_2(\varepsilon)-\varepsilon|\leqslant\frac{\varepsilon^2}{2},
\qquad
|\log(1-\varepsilon)+\varepsilon|\leqslant\varepsilon^2.
\]

第一式来自式 (8)，第二式来自对数幂级数的尾和估计。代入反射公式，所得余项 \(R(\varepsilon)\) 满足

\[
|R(\varepsilon)|\leqslant\varepsilon^2\left(|\log\varepsilon|+\frac12\right).
\]

这既证明原来的渐近式，也解释了为什么主修正项含 \(\varepsilon\log\varepsilon\)，不能假设函数在 \(x=1\) 处有通常的 Taylor 展开。

\hypertarget{ux8d1fux5b9eux8f74ux65e0ux7a77ux8fdcux5904}{%
\subsection{负实轴无穷远处}\label{ux8d1fux5b9eux8f74ux65e0ux7a77ux8fdcux5904}}

由式 (4) 与零点展开，当 \(a\to+\infty\) 时，

\[
\operatorname{Li}_2(-a)
=-\frac12\log^2a-\frac{\pi^2}{6}
+\frac1a-\frac1{4a^2}+O(a^{-3}).
\]

具体地，\(a\geqslant2\) 时，最后余项的绝对值不超过 \(1/[9a^3(1-1/a)]\)。因此原讲义的较粗形式

\[
\operatorname{Li}_2(-a)=-\frac12\log^2a-\frac{\pi^2}{6}+O(1/a)
\]

同样成立。这里是在小参数 \(-1/a\) 处使用收敛级数，再借反演公式返回 \(-a\)，并非在收敛区间外形式地展开。

\hypertarget{ux65b9ux6cd5ux603bux7ed3ux4e0eux53c2ux8003}{%
\section{方法总结与参考}\label{ux65b9ux6cd5ux603bux7ed3ux4e0eux53c2ux8003}}

特殊值计算的基本路径是：先寻找自变量之间的代数关系，再通过反射、Landen、倍角与反演公式建立方程组。求和时，分母中的 \(n\) 对应对生成函数作一次``除以自变量再积分''；带调和数的级数因此自然联系到二对数函数。

计算前先核对定义域，换序前先证明绝对收敛或非负性，取端点前先给出极限依据。若不能得到简单闭式，就把自变量约化到较小区间，并用式 (8) 控制截断误差。数值核对可以发现错误，但不能替代这些证明。

公式按本次文章的引用顺序统一编号。本次校订补全了实轴定义与所有上述推导，核心特殊值未改变。

参考：\href{https://dlmf.nist.gov/25.12\#i}{NIST DLMF §25.12(i)} 的式 25.12.1---25.12.6 可用于核对二对数的标准定义与连接公式。本文仅使用其对应的实轴情形，并在文内给出证明；进一步的级数方法见 \href{https://vincent-arqain.github.io/math-notes/notes/890e1516e64e21a2fb2248b4/}{Euler 型级数讲义}。

\end{document}
